% GENERATED FILE. Edit canonical lessons, then run npm run generate:books.
% canonical-source-hash: fnv1a64:54f84a3ac734f0f2
% canonical-lessons: UR-C54-chaval, UR-C54-dal, UR-C54-phal, UR-C54-seb, UR-C54-kela

\chapter{Rice, Lentils, Fruit, Apple, Banana}
\label{ch:ur-rice-lentils-fruit-apple-banana}
\hlchaptermodality{54}

\begin{chapteropening}
\textbf{By the end of this chapter:} \emph{I can say rice, lentils, fruit, an apple and a banana.}

Complete the last lesson of chapter 54: I can say rice, lentils, fruit, an apple and a banana.
\end{chapteropening}

\section[\texorpdfstring{\ur{چاول} (chāval)}{چاول (chāval)}]{\ur{چاول} (chāval) — rice}

\label{lesson:UR-C54-chaval}



\begin{quote}
Before the new one: say the Urdu for meat, then the Urdu for a fish.
\end{quote}

\subsection*{You'll want to know: \ur{چاول}}
\textbf{\ur{چاول}} — \emph{chāval} — \textquotedblleft{}rice\textquotedblright{}.

\begin{cousinweb}
\textbf{\ur{چاول}} (\emph{chāval}) is a tadbhava, a word inherited and reshaped by sound change. It comes from \emph{*cāvala}, \textquotedblleft{}husked rice\textquotedblright{}, a form CDIAL reconstructs from its descendants (CDIAL 4749). CDIAL suggests, with a \textquotedblleft{}possibly\textquotedblright{}, that it shares a non-Aryan origin with Sanskrit \emph{taṇḍula}, \textquotedblleft{}rice\textquotedblright{}.
\end{cousinweb}

\begin{grammarlens}[title={What you've built}]
One more word for this chapter.
\end{grammarlens}

\subsection*{Your turn}
\begin{itemize}
\raggedright
  \item \emph{Say it:} \emph{chāval}
  \item \emph{Say it:} \emph{chāval}, once more
  \item \emph{Say it:} say \emph{chāval}
  \item \emph{Recall:} say the Urdu for meat, then the Urdu for a fish, then say \emph{chāval} again
\end{itemize}

\begin{tcolorbox}[breakable,colback=teal!4,colframe=teal!35!black,title={Before you move on}]
What does \ur{چاول} mean? (\textquotedblleft{}Rice\textquotedblright{}.) Say it once more.
\end{tcolorbox}
\section[\texorpdfstring{\ur{دال} (dāl)}{دال (dāl)}]{\ur{دال} (dāl) — lentils}

\label{lesson:UR-C54-dal}



\begin{quote}
Before the new one: say the Urdu for a fish, then the Urdu for rice.
\end{quote}

\subsection*{You'll want to know: \ur{دال}}
\textbf{\ur{دال}} — \emph{dāl} — \textquotedblleft{}lentils\textquotedblright{}.

\begin{cousinweb}
\textbf{\ur{دال}} (\emph{dāl}) is a tadbhava, a word inherited from Sanskrit and reshaped by sound change. It comes from \emph{dāla}, \textquotedblleft{}kind of grain, split peas\textquotedblright{} (CDIAL 6309). CDIAL asks whether it is the same word as \emph{*dāla}, \textquotedblleft{}splitting\textquotedblright{}.
\end{cousinweb}

\begin{grammarlens}[title={What you've built}]
One more word for this chapter.
\end{grammarlens}

\subsection*{Your turn}
\begin{itemize}
\raggedright
  \item \emph{Say it:} \emph{dāl}
  \item \emph{Say it:} \emph{dāl}, once more
  \item \emph{Say it:} say \emph{dāl}
  \item \emph{Recall:} say the Urdu for a fish, then the Urdu for rice, then say \emph{dāl} again
\end{itemize}

\begin{tcolorbox}[breakable,colback=teal!4,colframe=teal!35!black,title={Before you move on}]
What does \ur{دال} mean? (\textquotedblleft{}Lentils\textquotedblright{}.) Say it once more.
\end{tcolorbox}
\section[\texorpdfstring{\ur{پھل} (phal)}{پھل (phal)}]{\ur{پھل} (phal) — fruit}

\label{lesson:UR-C54-phal}



\begin{quote}
Before the new one: say the Urdu for rice, then the Urdu for lentils.
\end{quote}

\subsection*{You'll want to know: \ur{پھل}}
\textbf{\ur{پھل}} — \emph{phal} — \textquotedblleft{}fruit\textquotedblright{}.

\begin{cousinweb}
\textbf{\ur{پھل}} (\emph{phal}) is a tadbhava, a word inherited from Sanskrit and reshaped by sound change. It comes from \emph{phala}, \textquotedblleft{}fruit\textquotedblright{} (CDIAL 9051). CDIAL compares Dravidian words for ripe fruit, such as Tamil \emph{paẓam}.
\end{cousinweb}

\begin{grammarlens}[title={What you've built}]
One more word for this chapter.
\end{grammarlens}

\subsection*{Your turn}
\begin{itemize}
\raggedright
  \item \emph{Say it:} \emph{phal}
  \item \emph{Say it:} \emph{phal}, once more
  \item \emph{Say it:} say \emph{phal}
  \item \emph{Recall:} say the Urdu for rice, then the Urdu for lentils, then say \emph{phal} again
\end{itemize}

\begin{tcolorbox}[breakable,colback=teal!4,colframe=teal!35!black,title={Before you move on}]
What does \ur{پھل} mean? (\textquotedblleft{}Fruit\textquotedblright{}.) Say it once more.
\end{tcolorbox}
\section[\texorpdfstring{\ur{سیب} (seb)}{سیب (seb)}]{\ur{سیب} (seb) — an apple}

\label{lesson:UR-C54-seb}



\begin{quote}
Before the new one: say the Urdu for lentils, then the Urdu for fruit.
\end{quote}

\subsection*{You'll want to know: \ur{سیب}}
\textbf{\ur{سیب}} — \emph{seb} — \textquotedblleft{}an apple\textquotedblright{}.

\begin{cousinweb}
\textbf{\ur{سیب}} (\emph{seb}) is a loan, usually traced to Persian \emph{sēb}, \textquotedblleft{}apple\textquotedblright{}.
\end{cousinweb}

\begin{grammarlens}[title={What you've built}]
One more word for this chapter.
\end{grammarlens}

\subsection*{Your turn}
\begin{itemize}
\raggedright
  \item \emph{Say it:} \emph{seb}
  \item \emph{Say it:} \emph{seb}, once more
  \item \emph{Say it:} say \emph{seb}
  \item \emph{Recall:} say the Urdu for lentils, then the Urdu for fruit, then say \emph{seb} again
\end{itemize}

\begin{tcolorbox}[breakable,colback=teal!4,colframe=teal!35!black,title={Before you move on}]
What does \ur{سیب} mean? (\textquotedblleft{}An apple\textquotedblright{}.) Say it once more.
\end{tcolorbox}
\section[\texorpdfstring{\ur{کیلا} (kelā)}{کیلا (kelā)}]{\ur{کیلا} (kelā) — a banana}

\label{lesson:UR-C54-kela}



\begin{quote}
Before the new one: say the Urdu for fruit, then the Urdu for an apple.
\end{quote}

\subsection*{You'll want to know: \ur{کیلا}}
\textbf{\ur{کیلا}} — \emph{kelā} — \textquotedblleft{}a banana\textquotedblright{}.

\begin{cousinweb}
\textbf{\ur{کیلا}} (\emph{kelā}) is a tadbhava, a word inherited from Sanskrit and reshaped by sound change. It comes from \emph{kadala}, \textquotedblleft{}banana plant\textquotedblright{} (CDIAL 2712). CDIAL takes the Sanskrit word from an Austroasiatic language.
\end{cousinweb}

\begin{grammarlens}[title={What you've built}]
One more word for this chapter.
\end{grammarlens}

\subsection*{Your turn}
\begin{itemize}
\raggedright
  \item \emph{Say it:} \emph{kelā}
  \item \emph{Say it:} \emph{kelā}, once more
  \item \emph{Say it:} say \emph{kelā}
  \item \emph{Recall:} say the Urdu for fruit, then the Urdu for an apple, then say \emph{kelā} again
\end{itemize}

\begin{tcolorbox}[breakable,colback=teal!4,colframe=teal!35!black,title={Before you move on}]
What does \ur{کیلا} mean? (\textquotedblleft{}A banana\textquotedblright{}.) Say it once more.
\end{tcolorbox}
